\documentclass[11pt,a4paper]{article}
\usepackage[margin=25mm]{geometry}
\usepackage{amsmath,amssymb,amsthm,fontspec,unicode-math,hyperref}
\setmainfont{STIXGeneral}
\setmathfont{STIX Two Math}
\setlength{\emergencystretch}{2em}
\newtheorem{theorem}{Theorem}
\newtheorem{lemma}[theorem]{Lemma}
\newcommand{\F}{\mathbb F_2}
\newcommand{\Aut}{\operatorname{Aut}}
\newcommand{\Cor}{\operatorname{Cor}}
\title{Graph groups and diagram symmetries}
\author{Human Compression research project}
\date{Working manuscript, 11 October 2026}
\begin{document}\maketitle
\begin{abstract}
The coordinate-kernel theorem for coset geometries applies directly to the
binary group attached to a simple graph. We use it to recover the diagram
realization of Leemans, Spiga and Tranchida. A residue-counting lemma gives
their order bound without a separate rank-three induction base. The isolated
minimum among all finite rank-three groups is reduced to three finite orders.
These arguments retain the original construction and Frucht's theorem.
\end{abstract}
\section{Sources and common theory}
Our principal source is Leemans--Spiga--Tranchida, \emph{Hypertopes with
prescribed diagram symmetries}, arXiv:2610.09711v1, nine pages. We recover
Theorems 1.1, 3.1--3.3, the displayed group order and the stated rank-three
minimum. The final assertion about proper quotients requires preservation
of the diagram; without that condition a Boolean quotient is a counterexample.
This qualification is not an answer to an open conjecture.

We use the complete common paper \emph{Supported kernels, translation lattices and finite marked quotients},
the companion common paper in this volume, especially UA.SK.Helly and UA.Geometry. The former
transfers the common-intersection property of typed cosets through a kernel
whose intersections with parabolics are coordinate subproducts. The latter
identifies the resulting geometry and its full type-preserving group. These
are existing project premises. Their full paper must be charged to a new
topic using this application, or once to Topic A when this application is
added there. The construction below is the same bilinear construction used
in UA.Forest; its coordinate part does not require a forest. No rigidity
claim for arbitrary quotients of a cyclic graph is inferred from UA.Forest.
We use right cosets throughout, as in that common paper and the source.

\section{The graph construction}
Let $X=(I,E)$ be a finite simple graph, with an ordering on $I$. Orient
an edge $e$ from its smaller endpoint $t(e)$ to its larger endpoint $h(e)$.
On $B_X=\F^I\times\F^E$ put
\[
 (v,z)(w,u)=(v+w,z+u+\beta(v,w)),\qquad
 \beta(v,w)_e=v_{h(e)}w_{t(e)}.
 \tag{1}
\]
\begin{lemma}
Equation (1) defines the group with generators $r_i$, and presentation
\[
 r_i^2=1,\qquad [[r_i,r_j],r_k]=1,\qquad
 [r_i,r_j]=1\quad(\{i,j\}\notin E).
 \tag{2}
\]
It has order $2^{|I|+|E|}$, and its subgroup on $J\subseteq I$ is exactly
\[
 B_J=\{(v,z):\operatorname{supp}(v)\subseteq J,
 \ \operatorname{supp}(z)\subseteq E[J]\}.
 \tag{3}
\]
\end{lemma}
\begin{proof}
Bilinearity gives
$\beta(v,w)+\beta(v+w,x)=\beta(w,x)+\beta(v,w+x)$, which is associativity.
The identity is $(0,0)$ and the inverse of $(v,z)$ is
$(v,z+\beta(v,v))$. Let $r_i=(e_i,0)$. These are involutions. Their
commutator on an edge is its coordinate vector in $\{0\}\times\F^E$;
other pairs commute. These commutators are central. Products of the $r_i$
in increasing order produce every $(v,0)$, and their commutators then
produce every central coordinate. Thus (2) maps onto this group.
In the abstract presentation its commutators are central. Conjugation by
$r_i$ both fixes $[r_i,r_j]$ and inverts it, so its square is one.
Collecting the generators in increasing order followed by edge commutators
therefore gives at most $2^{|I|+|E|}$ words. The coordinate group attains
this bound, proving the presentation and uniqueness. The same collection
and the coordinate generators internal to $J$ prove (3) in both directions.
The square of $(v,z)$ is $(0,\beta(v,v))$, so every element has fourth
power one. If $E$ is nonempty, an edge product has order four and a
nontrivial central commutator; the group then has exponent four and
nilpotency class exactly two.
\end{proof}

\begin{theorem}
The standard coset system of the marking $(r_i)_{i\in I}$ is a regular
hypertope $\Gamma_X$ with full type-preserving group $B_X$. Its diagram
is $X$ with all edges labelled four, and
\[
 \Cor(\Gamma_X)/\Aut(\Gamma_X)\cong\Aut(X).
 \tag{4}
\]
Every finite group occurs as this quotient for some connected $X$.
\end{theorem}
\begin{proof}
Formula (3) gives the intersection property. Project (1) to the vertex
coordinates $V=\F^I$. The images of the standard parabolics are coordinate
hyperplanes. Their typed cosets fix one coordinate each, so they have the
common-intersection property. The kernel is the product of the one-dimensional
edge groups, and its intersection with $B_J$ retains precisely $E[J]$.
The projection is onto every marked subgroup, and $V$ has the intersection
property. These verify all hypotheses of UA.SK.Helly. The markers are
nonidentity involutions, so UA.Geometry supplies thinness, residual
connectedness, regularity, and the asserted full group. The order of $r_ir_j$
is four on an edge, since its square is the nonzero edge coordinate, and
is two off an edge. This also identifies every rank-two residue and the diagram.

A graph automorphism preserves the complete presentation (2), hence permutes
the markers and the parabolics and induces a correlation of the coset geometry.
Conversely, compose any correlation with an element of the regular chamber
group so that it fixes the base chamber. Conjugation sends its unique
$i$-adjacency involution to the $\sigma(i)$-adjacency involution, where
$\sigma$ is the type permutation. Thus it preserves pair orders and induces
a graph automorphism. If $\sigma$ is trivial, the normalized correlation
fixes a chamber and is type preserving, so UA.Geometry makes it the identity.
These constructions are inverse on the quotient in (4).
Frucht's theorem supplies a finite connected simple graph with automorphism
group any prescribed finite group. If necessary the trivial group is represented
by an asymmetric connected tree with more than one vertex. This classical
realization theorem is retained, with its original attribution.
\end{proof}

\section{Growth in residues}
\begin{lemma}
Let a finite hypertope of rank $r\geq2$ have complete diagram, with every
rank-two polygon having at least $m\geq3$ elements of each type. Each type
has at least $m+r-2$ elements. If the hypertope is regular and its chamber
group is a two-group, each type has at least $2^r$ elements, including rank
one, and its group order is at least $2^{r(r+1)/2}$.
\end{lemma}
\begin{proof}
For the first assertion the rank-two base is the hypothesis. Suppose $r>2$
and fix a type $i$. Choose another type $j$ and a rank-two residue on $i,j$.
Such a residue exists by completing a flag. Its polygon has at least three
elements of each type, so it has a nonincident pair $x,y$ of types $i,j$.
The residue at $y$ has rank $r-1$ and inherits the same rank-two condition.
It contains at least $m+r-3$ elements of type $i$ by induction. The element
$x$ is outside that residue, which proves the asserted strict increase.

For the second assertion start in rank one with its two elements. The same
argument gives more than $2^{r-1}$ elements of each type at rank $r\geq2$:
use an $i,j$ residue with a nonincident pair, and apply the inductive bound
to the residue at $y$. This residue is regular under the stabilizer $G_y$ of $y$ in the original
chamber group. Indeed, $G_y$ is transitive on its chambers, and an element
in the action kernel fixes a full original chamber and is therefore trivial.
Thus its chamber group is a subgroup of the original two-group. Every type size in the whole regular
geometry is a subgroup index and hence a power of two. A power of two
strictly exceeding $2^{r-1}$ is at least $2^r$.

Choose a chamber and its successive subflags. Orbit--stabilizer identifies
successive stabilizer indices with type sizes in residues of ranks
$1,\ldots,r$. The full chamber stabilizer is trivial, and multiplication of
these indices gives the group order. The bounds just proved yield
$\prod_{k=1}^r2^k=2^{r(r+1)/2}$.
\end{proof}

For $X=K_r$, our group has exactly $2^{r+\binom r2}$ elements, so it
attains this bound among two-groups. More generally the same lower bound
holds for every regular two-group hypertope with a complete diagram; the
argument does not use the common label four. This is a direct extension
of the source's counting argument, not a claim of historical priority.

For the source's unrestricted rank-three minimum, take $m=4$. There are
at least five elements of every type, and each corresponding rank-two
residue has eight chambers. Thus the group order is at least forty and
is divisible by eight. Below sixty-four only orders $40,48,56$ remain.
The complete program accompanying this application examines all $79$ groups
of these orders in SmallGroups and all $32,628$ unordered triples of
nonidentity involutions. There are $88$ triples with all pair products of
order four, all in groups of order forty-eight; none generates its group.
Every regular rank-three hypertope would supply such a generating triple.
Consequently no smaller group occurs. This finite statement depends on
the SmallGroups classification and GAP's exact group algorithms. It does
not use the source's external hypermap list, and is separate from the
uniform written lower bound.

Finally a proper marked quotient retaining the complete all-four diagram
cannot be a hypertope, by the two-group bound. Without preservation of that
diagram the conclusion fails: $B_{K_r}\to\F^r$ kills all edge coordinates.
The Boolean maximal-parabolic cosets form the join of $r$ pairs of points,
a regular hypertope with empty diagram. The quotient is proper for $r\geq2$.
This gives the exact scope of the closing source assertion.
\section*{Attribution and verification}
The graph group, its presentation, the diagram realization and the two-group
minimum belong to Leemans, Spiga and Tranchida. The source credits Li and
Qiao for the construction's inspiration. Here the existing coordinate-kernel
argument supplies flag transitivity in one application, and residue growth
removes the special induction base. Small-order exclusion is independently
recomputed. The companion Lean file verifies the bilinear group law over
arbitrary additive abelian groups, its vertex projection, the binary edge
form, the support subgroups and their intersection law, and the numerical
power-of-two step. It does not identify those subgroups with the generated
parabolics or formalize the presentation, geometry, correlation theorem or
finite search. No complete formalization, independent review, manuscript
admission or measured compression factor is assigned to this working draft.
\begin{thebibliography}{9}
\bibitem{LST} D. Leemans, P. Spiga and P. Tranchida, Hypertopes with prescribed
diagram symmetries, arXiv:2610.09711v1 (2026).
\bibitem{UA} Human Compression research project, Supported kernels, translation lattices and finite marked quotients, complete companion common paper in this volume (2026).
\bibitem{Frucht} R. Frucht, Herstellung von Graphen mit vorgegebener abstrakter
Gruppe, Compositio Mathematica 6 (1939), 239--250.
\bibitem{LQ} Y. Li and Y. Qiao, Group-theoretic generalisations of vertex and
edge connectivities, Proceedings of the AMS 148 (2020), 4679--4693.
\end{thebibliography}
\end{document}
